\scanpage{269}
\mspage{137}

\diagram{\makebox[0pt][r]{$(159)$\qquad}\begin{tikzcd}[ampersand replacement=\&,row sep=small,column sep=small]
\& \unsure{e|\mathrm e}\arrow[r,hook,no head]\arrow[d,hook,no head] \& U_0\arrow[d,hook,no head]\\
\& T_H\arrow[r,hook,no head]\arrow[d,hook,no head] \& H\arrow[d,hook,no head]\\
T_0\arrow[r,hook,no head] \& T_G\arrow[r,hook,no head] \& G
\end{tikzcd}\ .}

Bien sûr, $\sigma$ n'est pas déterminé en termes des données précédentes, mais on a

\[
\left\{\begin{aligned}
N_{\deleted{H}G}&=\{g\in\deleted{H}G\mid u(\{D_0,D_1\})=\{D_0,D_1\}\},\quad\text{d'où}\\
&\deleted{\text{donc }N_H\unsure{=}}\\
N_H&=N_G\cap H=\{g\in H\mid u(\{D_0,D_1\})=\{D_0,D_1\}
\end{aligned}\right.\tag*{(160)}
\]

l'hom

\[
N_G\longrightarrow\{\pm1\}\qquad
\begin{aligned}&(\text{si }\sigma\notin H\unsure{\}},\text{ i.e.}\\&D_0\ne D_1)\end{aligned}
\]

étant donné par l'action de $g$ sur la paire $D_0,D_1$ \ldots. Ainsi $N_G$, $N_H$ et l'hom $N_G\to\pm1$ \deleted{de} de noyau $T_G$ (resp. $N_H\to\pm1$ de noyau $T_H$) sont déterminés par les données. Ceci dit, on a des \deleted{bijections}\inserted{isom} canoniques de $Z$-torseurs

\[
\begin{aligned}
&\deleted{\unsure{D_1\xrightarrow{\sim}\Pi}\ \unsure{\xrightarrow{\sim}}}\quad
\underset{\substack{\deleted{\unsure{N_H\setminus H}}\\\text{image inverse de }-1\\\text{par }N_H\to\{\pm1\}}}{\underbrace{N_H^-}}\xrightarrow{\sim}D_1\xrightarrow{\sim}\unsure{\Pi|\pi}\\
&\deleted{t_0\mapsto\unsure{\tau}(s_0,t_0)}\hspace{5em}
\sigma\longmapsto\sigma(s_0),\ t_0\longmapsto\tau(s_0,t_0)
\end{aligned}\tag*{(\unsure{151|161})}
\]

\deleted{l'isom} \deleted{obtenu par} où on prend sur $N_H^-$ la structure de $Z$ torseur \unsure{défini|définie} par les opérations \uline{droites} de $T_H\xleftarrow[u]{\sim}Z$ sur $N_H^-$, i.e. \deleted{\unsure{$\lambda$}}

\[
\lambda\in Z\ \text{opère sur }N_H^-\ \text{par}\quad g\longmapsto gu(\lambda).\tag*{(162)}
\]

Notons ici :
% archive sheet number: 268
